\section{Scheduling the first positive answer}
\label{sec:scheduling}

The maximum-depth algorithm schedules discovery around the heights
of its child codes. For expected depth, a cap containing a likely
target should instead be encountered early. This is a weighted min-sum set cover (MSSC) problem. The sets remain
the supplied caps;
singleton sets are not added unless an actual supplied test realizes
them.

In the decision-tree application, a positive hull answer exits the
all-negative route and reveals that the target lies in a small cap.
The all-zero pattern has no positive answer, so on the full-cover
route it is identified after every positive pattern has been covered
(and an earlier singleton inference only shortens that route). Its
contribution is therefore at most its weight times the schedule
length. This is why the
ordering analysis keeps both the first-cover sum and a separate
length term.

\begin{symboltable}
$U,\ \pi_h,\ p_0$ & Coverable positive-pattern demands, their positive weights, and the weight of a target found only after the schedule finishes.\\
$t_Q(h),\ \MSSC(Q)$ & First-cover position of demand $h$ in ordering $Q$ and the sum $\sum_h\pi_h t_Q(h)$.\\
$k_0$ & Minimum number of candidate tests needed to cover all demands in $U$.\\
$\beta,\ \Phi_\beta,\ Z_\beta$ & Geometric budget factor and the first-cover and total-length accounting functions.\\
$r_Q(x),\ G_u$ & Comparison-order mass with first-cover time greater than $x$, and the constructed schedule for shift $u$.\\
\end{symboltable}

For background, weighted MSSC has a classical greedy
four-approximation~\cite{flt2004}. We do not use that guarantee in
either proof. Instead, exact coverage with a cardinality budget
gives the stronger circular-interval ordering proved below,
including a bound on its total length.

\begin{lemma}[Exact budgeted interval coverage]\label{lem:budget-coverage}
For positively weighted interval or circular-interval demands on
a finite candidate set, a set of at most $k$ candidates hitting
maximum demand weight is computable in polynomial time.
\end{lemma}
\begin{proof}
Algorithm~\ref{alg:interval-cover} gives both interval oracles.
For ordinary intervals, number the candidates $1,\ldots,M$.
If the previous selected candidate is $p$ and the next is $q>p$,
the newly hit demands are exactly the intervals containing $q$
whose left endpoint is greater than $p$. An interval containing
$q$ and an earlier selected point must also contain $p$.
Conversely, any interval containing $q$ with left endpoint greater
than $p$ contains no earlier selected candidate. Thus the added
weight $g(p,q)$ depends only on these two positions, not on the
earlier choices.

For exactly $t$ selected candidates ending at $q$, use the recurrence
\[
 D(t,q)=\max_{p<q}\{D(t-1,p)+g(p,q)\},
 \qquad D(0,0)=0.
\]
The position $0$ is a \emph{sentinel}, an artificial initial marker
before every interval, not a selectable candidate.
All other zero-candidate or nonexistent states are infeasible
and have value $-\infty$ in this maximization.
Take the maximum over
$t\le k$ and final positions, including the empty selection.
This is a polynomial dynamic program.

For circular intervals and $k\ge1$, enumerate a selected candidate
$v$. Credit the demands it hits and delete them. Cutting at $v$
turns the remaining demands into ordinary intervals, on which
the dynamic program selects at most $k-1$ more candidates.
One enumeration contains a candidate of an optimal nonempty
selection, so taking the best is exact. The $k=0$ case is empty.
\end{proof}

\begin{algorithm}[H]
\caption{\IntervalCover: exact minimum cover or budgeted coverage}
\label{alg:interval-cover}
\KwIn{Ordered interval demands with weights; either \emph{minimum} or budget $k$;
ordinary or circular order.}
\KwOut{A minimum hitting set, or a maximum-weight set of size at most $k$.}
\If{order is circular}{
  Enumerate a forced selected candidate $v$; delete its hit demands and cut there\;
  Solve each residual ordinary instance, with budget $k-1$ in budget mode\;
  \KwRet the best residual solution plus $v$, also considering the empty budget solution\;
}
\eIf{mode is \emph{minimum}}{
  \If{an empty demand exists}{\KwRet infeasible\;}
  $P\gets\varnothing$\;
  \While{a demand remains}{
    Choose the right endpoint $q$ of an earliest-ending remaining demand\;
    $P\gets P\cup\{q\}$; delete all demands containing $q$\;
  }
  \KwRet $P$\;
}{
  Set $D(0,0)=0$ and all other unreachable states to $-\infty$\;
  Precompute $g(p,q)$: weight of intervals containing $q$ with left endpoint $>p$\;
  \For{$t=1,\ldots,\min\{k,M\}$}{
    \For{$q=1,\ldots,M$}{
      $D(t,q)\gets\max_{p<q}(D(t-1,p)+g(p,q))$; store a maximizing $p$\;
    }
  }
  \KwRet the reconstruction of the best state with $t\le k$, including $(0,0)$\;
}
\end{algorithm}

For minimum circular cover the forced candidate need only range over
one nonempty demand. For budget coverage, empty demands contribute
zero covered weight. The empty family returns the empty set in both
modes. Lemmas~\ref{lem:minimum-interval-cover} and
\ref{lem:budget-coverage} establish exactness.


The next lemma uses exact budgeted covers at geometrically
increasing budgets. Small budgets can discover likely targets early;
the minimum cover at the final budget guarantees that every positive
target is eventually found. Shifting the budget sequence changes
where it crosses integer thresholds. Averaging over shifts bounds
both positive first-cover times and the zero-pattern waiting time;
a finite search over the distinct resulting schedules removes the
random shift.

\begin{lemma}[Geometric-batch ordering]\label{lem:mssc}
Let $U$ be a nonempty family of coverable weighted circular intervals.
Let $k_0$ be its minimum cover size and let $p_0\ge0$ be the weight
of a target covered only when the entire schedule finishes.
Fix rational $\beta>1$, and put
\[
 r_\beta=\lfloor\beta\rfloor,\qquad
 \lambda_\beta=\frac{\beta+1}{2\ln\beta},\qquad
 \nu_\beta=1+\frac1{\ln\beta},\qquad
 \eta_\beta=\frac{r_\beta+1}{2}
            -\frac{\ln(r_\beta!)}{2\ln\beta}-\lambda_\beta.
\]
For integer $t,k\ge1$, define the accounting functions
\[
 \begin{aligned}
 \Phi_\beta(t)
 &=\lambda_\beta+\eta_\beta
    +\frac1{2\ln\beta}\int_1^t
       \frac{\lfloor x\rfloor+\lfloor\beta x\rfloor}{x}\,dx,\\
 Z_\beta(k)
 &=k+\frac{k\ln k-\ln(k!)}{\ln\beta}.
 \end{aligned}
\]
A deterministic polynomial-time ordering $G$ satisfies, for every
comparison cover ordering $Q$,
\[
 \MSSC(G)+p_0|G|
 \le\sum_{h\in U}\pi_h\Phi_\beta(t_Q(h))+p_0 Z_\beta(k_0).
\]
The simpler bounds
$\Phi_\beta(t)\le\lambda_\beta t+\eta_\beta$ and
$Z_\beta(k)\le\nu_\beta k-1/\ln\beta$ also hold.
Every candidate schedule considered by the construction also has
length at most $\bigl(1+\beta/(\beta-1)\bigr)k_0$.
\end{lemma}
\begin{proof}
Compute a minimum cover $Q_*$ using
Lemma~\ref{lem:minimum-interval-cover} and, for each integer $k<k_0$,
an exact maximum-weight coverage set $Q_k$ of at most $k$
candidates using Lemma~\ref{lem:budget-coverage}.
Fix the selected sets deterministically.

Fix a comparison cover ordering $Q$, and write
\[
 r_Q(x)=\sum_{h:t_Q(h)>x}\pi_h
\]
for the weight whose first-cover index in $Q$ exceeds $x$.
For a shift $u\in[1,\beta)$, use budgets
$k^{(j)}=\lfloor\beta^ju\rfloor$, $j=0,1,\ldots$.
In each phase, process $Q_{k^{(j)}}$ when $k^{(j)}<k_0$ and $Q_*$
otherwise. Order that set greedily by the weight it covers among
the demands still unhit by earlier queries. Omit repeated queries
and those hitting no remaining demand. Stop after the first phase
whose budget reaches $k_0$.
\begin{claim}[Phase coverage and schedule length]
After a phase with budget $k<k_0$, the weight still unhit is at
most $r_Q(k)$. The complete shifted schedule has length at most
$\bigl(1+\beta/(\beta-1)\bigr)k_0$.
\end{claim}
\begin{proof}
The union of all queries processed through this phase covers every
demand covered by $Q_k$: a query omitted for hitting no remaining
demand has its covered demands already hit by an earlier phase.
The set $Q_k$ covers at least as much total weight as the first
$k$ queries of $Q$, since those queries are one feasible
at-most-$k$ candidate set. Hence the schedule has hit at least the
weight hit by that comparison prefix, leaving at most $r_Q(k)$.
The final $Q_*$ phase covers all demands.

Let $J$ be the first phase whose budget reaches $k_0$. The
pre-final phases use at most
\[
 \sum_{j<J}\lfloor\beta^ju\rfloor
 <\sum_{j<J}\beta^ju
 <\frac{\beta}{\beta-1}k_0
\]
queries: the last pre-final unrounded budget is below $k_0$, and
the preceding terms decrease geometrically by a factor $1/\beta$.
The final minimum cover uses $k_0$ queries, proving the length
bound.
\end{proof}

\begin{claim}[First-cover cost within one batch]
Suppose a phase processes $m$ queries, newly hitting weights
$g_1,\ldots,g_m$ in order, and starts and ends with unhit weights
$r$ and $r'$. Its contribution to the first-cover sum is at most
\[
 mr'+\sum_{i=1}^m i g_i
 \le mr'+\frac{m+1}{2}(r-r')
 =\frac{m+1}{2}r+\frac{m-1}{2}r'.
\]
\end{claim}
\begin{proof}
After each greedy choice, the remaining coverage of every unused
candidate can only decrease. Therefore the newly hit weights are
nonincreasing: $g_1\ge\cdots\ge g_m$. Pairing ranks $i$ and
$m+1-i$ shows
\[
 \sum_{i=1}^m i g_i\le\frac{m+1}{2}\sum_{i=1}^m g_i.
\]
Indeed, in each pair the earlier rank has at least as much weight
as the later one, so replacing both ranks by their mean cannot
decrease the weighted sum. A demand first hit at rank $i$ remains
unhit for exactly $i$ queries in this batch, while each demand
still unhit at the end remains unhit for all $m$ queries. Since
$\sum_i g_i=r-r'$, their combined contribution gives the claim.
\end{proof}

Now let $r_j$ be the unhit weight after phase $j$, and put
$r_{-1}=w(U)$. If phase $j$ processes $m_j$ queries, the preceding
claim bounds its first-cover contribution by
\[
 \frac{m_j+1}{2}r_{j-1}+\frac{m_j-1}{2}r_j.
\]
Since $m_j\le k^{(j)}$, replacing $m_j$ by its integer budget and
combining adjacent phase coefficients gives
\[
 \MSSC(G_u)\le
 \frac{\lfloor u\rfloor+1}{2}w(U)
 +\frac12\sum_{j\ge0}
       \bigl(k^{(j)}+k^{(j+1)}\bigr)r_Q(\beta^ju).
\]
For every pre-final phase, the first claim gives
$r_j\le r_Q(k^{(j)})=r_Q(\beta^ju)$; the equality uses the
integer-valued comparison first-cover indices. The final phase
has residual zero. The displayed infinite sum only adds
nonnegative upper bounds for the final and later phases.

\paragraph{Average the first-cover cost.}
Average this expression over a \emph{log-uniform shift}, with
probability density $1/(u\ln\beta)$ on $[1,\beta)$. This density
integrates to one because $\int_1^\beta du/u=\ln\beta$.
Splitting the first term's integral at its integers
gives
\[
 \begin{aligned}
 \mathbb E[\lfloor u\rfloor]
 &=\frac{\displaystyle
     \sum_{k=1}^{r_\beta-1}k\ln\frac{k+1}{k}
        +r_\beta\ln\frac{\beta}{r_\beta}}{\ln\beta}\\
 &=r_\beta-\frac{\ln(r_\beta!)}{\ln\beta}.
 \end{aligned}
\]
The intervals $[\beta^j,\beta^{j+1})$ partition $[1,\infty)$,
so averaging the remaining terms gives
\[
 \begin{aligned}
 \mathbb E[\MSSC(G_u)]
 &\le(\lambda_\beta+\eta_\beta)w(U)\\
 &\quad+\frac1{2\ln\beta}\int_1^\infty
      \frac{\lfloor x\rfloor+\lfloor\beta x\rfloor}{x}
      r_Q(x)\,dx\\
 &=\sum_{h\in U}\pi_h\Phi_\beta(t_Q(h)).
 \end{aligned}
\]
For each real $x$, write
$r_Q(x)=\sum_{h:t_Q(h)>x}\pi_h$. Interchanging this finite sum
with the integral makes each hypothesis contribute the integral
from $1$ to its integer first-cover time. Adding the initial
$(\lambda_\beta+\eta_\beta)w(U)$ term gives exactly
$\sum_{h\in U}\pi_h\Phi_\beta(t_Q(h))$.
Bounding the integrand by $\beta+1$ recovers
$\Phi_\beta(t)\le\lambda_\beta t+\eta_\beta$, but would discard
the savings from all later integer budgets.

\begin{claim}[Shift-averaged schedule length]
Under the same shift distribution,
\[
 \mathbb E[|G_u|]\le Z_\beta(k_0)
 \le\nu_\beta k_0-\frac1{\ln\beta}.
\]
\end{claim}
\begin{proof}
The final cover contributes $k_0$ queries. The earlier phases
contribute at most
$\sum_{j:\,\beta^ju<k_0}\lfloor\beta^ju\rfloor$.
After averaging and substituting $x=\beta^ju$, their expected
length is
\[
 \frac1{\ln\beta}\int_1^{k_0}\frac{\lfloor x\rfloor}{x}\,dx.
\]
On each interval $[j,j+1)$, $\lfloor x\rfloor=j$, so summing the
integrals gives
\[
 \int_1^{k_0}\frac{\lfloor x\rfloor}{x}\,dx
 =k_0\ln k_0-\ln(k_0!).
\]
Consequently $\mathbb E[|G_u|]\le Z_\beta(k_0)$. Since $\ln x$
is increasing,
\[
 \ln(k_0!)=\sum_{j=2}^{k_0}\ln j
 \ge\int_1^{k_0}\ln x\,dx
 =k_0\ln k_0-k_0+1.
\]
Substituting this in $Z_\beta(k_0)$ gives
$Z_\beta(k_0)\le\nu_\beta k_0-1/\ln\beta$. The formula also gives
$Z_\beta(1)=1$.
\end{proof}

The first-cover and length bounds hold under the same shift
distribution. Therefore
\[
 \mathbb E[\MSSC(G_u)+p_0|G_u|]
 \le\sum_{h\in U}\pi_h\Phi_\beta(t_Q(h))
      +p_0Z_\beta(k_0).
\]

No random shift is used by the algorithm. Let $J$ be the first
integer with $\beta^J\ge k_0$. All budget changes occur at the
rational breakpoints $u=k/\beta^j$ for
$0\le j\le J$, $1\le k\le k_0$, inside $[1,\beta)$.
Enumerate these breakpoints and one rational midpoint between
each consecutive pair, construct the resulting orderings, and
take the one minimizing the rational quantity
$\MSSC(G_u)+p_0|G_u|$. There are
$O(k_0\log k_0)$ intervals for fixed $\beta$.
The budgets, selected sets, and their greedy internal orders are
constant on each interval. Their minimum joint value is at most
the average just bounded. The waiting-time weight $p_0$ is used
only to choose the schedule; all queries remain supplied candidates.
\end{proof}

\begin{algorithm}[H]
\caption{\BatchOrders: deterministic family of geometric-batch schedules}
\label{alg:batch-orders}
\KwIn{Weighted circular demands and rational $\beta>1$.}
\KwOut{All schedules needed for deterministic shift minimization.}
Compute a minimum cover $Q_*$ of size $k_0$ using $\IntervalCover$\;
\For{$k=1,\ldots,k_0-1$}{Compute a maximum-weight budget-$k$ set $Q_k$\;}
Find the first $J$ with $\beta^J\ge k_0$ by rational multiplication\;
Form the breakpoints $1,\beta$ and $k/\beta^j\in[1,\beta)$,
$0\le j\le J$, $1\le k\le k_0$\;
\ForEach{breakpoint below $\beta$ and each midpoint of consecutive breakpoints}{
  $u\gets$ that rational representative; $Q\gets$ empty order\;
  \For{$j=0,1,\ldots$}{
    $k\gets\lfloor\beta^ju\rfloor$;
    $P\gets Q_k$ if $k<k_0$, and $P\gets Q_*$ otherwise\;
    Greedily order $P$ by newly covered weight, append its productive unused candidates to $Q$\;
    \If{$k\ge k_0$}{stop this shift's phases\;}
  }
  Retain $Q$\;
}
\KwRet the retained family, removing duplicate schedules\;
\end{algorithm}

An empty demand family needs no schedule. Lemma~\ref{lem:mssc}
bounds the family's minimum joint first-cover and waiting cost.
The number of representatives is $O(k_0\log k_0)$ for fixed $\beta$.

